\documentclass[10pt]{article}
%Formato extenso:report
\begin{sloppypar}
	
\end{sloppypar}
%Formato corto:article

% Esto es para que el LaTeX sepa que el texto esten espanol:
\usepackage[spanish]{babel}

% Esto es para poder escribir acentos directamente, no recomendable si trabaja en varios sistemas operativos:
%\usepackage[ansinew]{inputenc}    % Windows: Descomentar si quiere escribir con acentos directamente
% \usepackage[latin1]{inputenc}         % Linux: Descomentar si quiere escribir con acentos directamente
%\usepackage[applemac]{inputenc} % MacOS: Descomentar si quiere escribir con acentos directamente

% Paquetes de la AMS:
\usepackage{amsmath, amsthm, amsfonts,amssymb}

% Borrame si quieres:
\usepackage{multicol}

% Referencias
\usepackage{hyperref}

% Paquete para escribir codigo
\usepackage{listings}
\lstset{basicstyle=\footnotesize\ttfamily,breaklines=true}
\usepackage{alltt}

% Paquete para incluir imagenes
\usepackage{graphicx}

% Paquete para incluir varias imagenes en una
% \usepackage{subfig}#subfig deprecated?
\usepackage{caption}
\usepackage{subcaption}

% para poder fijare las imagenes ([H])
\usepackage{float}

% para agregar opciones al enumerate
\usepackage{enumerate}


% Teoremas
\newtheorem{thm}{Teorema}[section]
\newtheorem{cor}[thm]{Corolario}
\newtheorem{lem}[thm]{Lema}
\newtheorem{prop}[thm]{Proposici\'on}
\theoremstyle{definition}
\newtheorem{defn}[thm]{Definici\'on}
\theoremstyle{remark}
\newtheorem{rem}[thm]{Observaci\'on}
\theoremstyle{definition}
\newtheorem{prob}{Problema}

% Calculus symbols
\newcommand{\pd}[2]{\frac{\partial #1}{ \partial #2}}   % First partial derivative command
\newcommand{\td}[2]{\frac{\mathrm{d} #1}{ \mathrm{d} #2}}
\newcommand{\pdd}[2]{\frac{\partial^2 #1}{ \partial #2 ^2}}   % Second partial derivative command
\newcommand{\pddc}[3]{\frac{\partial^2 #1}{ \partial #2 \partial #3}}   % Second partial derivative command

% Continuum mechanics & FEM symbols
\def\sca   #1{\mbox{\rm{#1}}{}}
\def\mat   #1{\mbox{\boldmath $\mathsf #1$}}
\def\vec   #1{\mbox{\boldmath $#1$}{}}
\def\ten   #1{\mbox{\boldmath $#1$}{}}
\def\ltr   #1{\mbox{\sf{#1}}}
\def\bltr  #1{\mbox{\sffamily{\bfseries{{#1}}}}}

% math operators and symbols
\DeclareMathOperator{\dive}{div}
\DeclareMathOperator{\trace}{trace}
\DeclareMathOperator{\tr}{tr}
\DeclareMathOperator{\symm}{symm}
\DeclareMathOperator{\sk}{skew}
\DeclareMathOperator{\grad}{grad}
\DeclareMathOperator{\Grad}{Grad}
\DeclareMathOperator{\curl}{curl}
\DeclareMathOperator{\Curl}{Curl}
\def\R{\mbox{\(\mathbb{R}\)}}
\def\dx{\mbox{\(\,\mathrm{d}x\)}}



% Page margins
%\topmargin=-2cm
%\textheight=22cm
%\textwidth=17cm
%\oddsidemargin=0cm
\usepackage{geometry}
\geometry{left=2.5cm, right=2.5cm, top=2cm, bottom=3cm}

\begin{document}
%   \input{template.tex}
  \begin{titlepage}
    \begin{figure}
      \begin{minipage}{2.5cm}
	\includegraphics[width=0.8\textwidth]{./figures/LogoUC-BN}
      \end{minipage}
      \begin{minipage}{14.5cm}
	\vspace{4mm}
	{\sc PONTIFICIA UNIVERSIDAD CAT\'OLICA DE CHILE}\\
	Escuela de Ingenier\'ia\\
	Departamento de Ingenier\'ia Estructural y Geot\'ecnica\\
	{\bf ICE 3233 Elementos Finitos}\\
	\vspace{0mm}
	\hrulefill
      \end{minipage}
    \end{figure}
    \phantom{""}\vspace{60mm}

    \begin{center}
      \Huge{\textbf{Tarea 8}}\vspace{95mm}\\
      \raggedleft \Large{Jos\'e Galaz Mora}\\ \today\\ Tiempo dedicado: 16 horas
    \end{center}
  \end{titlepage}
  


  \input{p1.tex}
  \input{p2.tex}
  \pagebreak

  \twocolumn
  \begin{center}
	  {\bf C\'odigos de Python}
  \end{center}

  \lstinputlisting[language=Python]{./codes/p1.py}
% %   \pagebreak
  \lstinputlisting[language=Python]{./codes/p2.py}
%   \pagebreak
  \lstinputlisting[language=Python]{./codes/tri3poisson.py}
%   \pagebreak
  \lstinputlisting[language=Python]{./codes/t8.py}
%   \pagebreak
%   \lstinputlisting[language=Python]{./codes/func_t6.py}
  
%   \lstinputlisting[language=Python]{./codes/func.py}
  
\end{document}